\begin{abstract}
\noindent
We study the causal effect of digital information provision on household electricity consumption using a randomized field experiment with a staggered rollout. Households were randomly assigned to an information condition (\(N = 449\)), receiving access to a smartphone application with feedback on electricity use, costs, emissions, social comparisons, and gamified learning features, or to a control condition (\(N = 443\)) without app access. We measure electricity consumption using high-frequency smart-meter data for up to 70 weeks after treatment. App access reduces daily electricity consumption by approximately 4.5--4.7\%, with reductions concentrated in daytime hours (07:00--22:00), when households have greater scope to adjust consumption. Point estimates remain negative throughout the observation period, but statistical significance is concentrated in the first six months after rollout. App engagement declines sharply after rollout, and higher engagement is associated with larger and more persistent reductions in consumption. The findings suggest that digital feedback can reduce residential electricity demand, but that persistence depends on sustained user attention.


\end{abstract}

\medskip
\noindent
\textbf{JEL Codes:} C93, D12, Q41 \\
\textbf{Keywords:} Energy consumption, Information provision, Energy conservation, Information salience, Environmental behavior, Randomized controlled trials